% Options for packages loaded elsewhere
\PassOptionsToPackage{unicode}{hyperref}
\PassOptionsToPackage{hyphens}{url}
%
\documentclass[
  10pt,
]{article}
\usepackage{amsmath,amssymb}
\usepackage{iftex}
\ifPDFTeX
  \usepackage[T1]{fontenc}
  \usepackage[utf8]{inputenc}
  \usepackage{textcomp} % provide euro and other symbols
\else % if luatex or xetex
  \usepackage{unicode-math} % this also loads fontspec
  \defaultfontfeatures{Scale=MatchLowercase}
  \defaultfontfeatures[\rmfamily]{Ligatures=TeX,Scale=1}
\fi
\usepackage{lmodern}
\ifPDFTeX\else
  % xetex/luatex font selection
  \setmainfont[]{DejaVu Serif}
  \setmonofont[]{DejaVu Sans Mono}
\fi
% Use upquote if available, for straight quotes in verbatim environments
\IfFileExists{upquote.sty}{\usepackage{upquote}}{}
\IfFileExists{microtype.sty}{% use microtype if available
  \usepackage[]{microtype}
  \UseMicrotypeSet[protrusion]{basicmath} % disable protrusion for tt fonts
}{}
\makeatletter
\@ifundefined{KOMAClassName}{% if non-KOMA class
  \IfFileExists{parskip.sty}{%
    \usepackage{parskip}
  }{% else
    \setlength{\parindent}{0pt}
    \setlength{\parskip}{6pt plus 2pt minus 1pt}}
}{% if KOMA class
  \KOMAoptions{parskip=half}}
\makeatother
\usepackage{xcolor}
\usepackage[margin=19mm]{geometry}
\usepackage{color}
\usepackage{fancyvrb}
\newcommand{\VerbBar}{|}
\newcommand{\VERB}{\Verb[commandchars=\\\{\}]}
\DefineVerbatimEnvironment{Highlighting}{Verbatim}{commandchars=\\\{\}}
% Add ',fontsize=\small' for more characters per line
\newenvironment{Shaded}{}{}
\newcommand{\AlertTok}[1]{\textcolor[rgb]{1.00,0.00,0.00}{\textbf{#1}}}
\newcommand{\AnnotationTok}[1]{\textcolor[rgb]{0.38,0.63,0.69}{\textbf{\textit{#1}}}}
\newcommand{\AttributeTok}[1]{\textcolor[rgb]{0.49,0.56,0.16}{#1}}
\newcommand{\BaseNTok}[1]{\textcolor[rgb]{0.25,0.63,0.44}{#1}}
\newcommand{\BuiltInTok}[1]{\textcolor[rgb]{0.00,0.50,0.00}{#1}}
\newcommand{\CharTok}[1]{\textcolor[rgb]{0.25,0.44,0.63}{#1}}
\newcommand{\CommentTok}[1]{\textcolor[rgb]{0.38,0.63,0.69}{\textit{#1}}}
\newcommand{\CommentVarTok}[1]{\textcolor[rgb]{0.38,0.63,0.69}{\textbf{\textit{#1}}}}
\newcommand{\ConstantTok}[1]{\textcolor[rgb]{0.53,0.00,0.00}{#1}}
\newcommand{\ControlFlowTok}[1]{\textcolor[rgb]{0.00,0.44,0.13}{\textbf{#1}}}
\newcommand{\DataTypeTok}[1]{\textcolor[rgb]{0.56,0.13,0.00}{#1}}
\newcommand{\DecValTok}[1]{\textcolor[rgb]{0.25,0.63,0.44}{#1}}
\newcommand{\DocumentationTok}[1]{\textcolor[rgb]{0.73,0.13,0.13}{\textit{#1}}}
\newcommand{\ErrorTok}[1]{\textcolor[rgb]{1.00,0.00,0.00}{\textbf{#1}}}
\newcommand{\ExtensionTok}[1]{#1}
\newcommand{\FloatTok}[1]{\textcolor[rgb]{0.25,0.63,0.44}{#1}}
\newcommand{\FunctionTok}[1]{\textcolor[rgb]{0.02,0.16,0.49}{#1}}
\newcommand{\ImportTok}[1]{\textcolor[rgb]{0.00,0.50,0.00}{\textbf{#1}}}
\newcommand{\InformationTok}[1]{\textcolor[rgb]{0.38,0.63,0.69}{\textbf{\textit{#1}}}}
\newcommand{\KeywordTok}[1]{\textcolor[rgb]{0.00,0.44,0.13}{\textbf{#1}}}
\newcommand{\NormalTok}[1]{#1}
\newcommand{\OperatorTok}[1]{\textcolor[rgb]{0.40,0.40,0.40}{#1}}
\newcommand{\OtherTok}[1]{\textcolor[rgb]{0.00,0.44,0.13}{#1}}
\newcommand{\PreprocessorTok}[1]{\textcolor[rgb]{0.74,0.48,0.00}{#1}}
\newcommand{\RegionMarkerTok}[1]{#1}
\newcommand{\SpecialCharTok}[1]{\textcolor[rgb]{0.25,0.44,0.63}{#1}}
\newcommand{\SpecialStringTok}[1]{\textcolor[rgb]{0.73,0.40,0.53}{#1}}
\newcommand{\StringTok}[1]{\textcolor[rgb]{0.25,0.44,0.63}{#1}}
\newcommand{\VariableTok}[1]{\textcolor[rgb]{0.10,0.09,0.49}{#1}}
\newcommand{\VerbatimStringTok}[1]{\textcolor[rgb]{0.25,0.44,0.63}{#1}}
\newcommand{\WarningTok}[1]{\textcolor[rgb]{0.38,0.63,0.69}{\textbf{\textit{#1}}}}
\setlength{\emergencystretch}{3em} % prevent overfull lines
\providecommand{\tightlist}{%
  \setlength{\itemsep}{0pt}\setlength{\parskip}{0pt}}
\setcounter{secnumdepth}{-\maxdimen} % remove section numbering
\ifLuaTeX
  \usepackage{selnolig}  % disable illegal ligatures
\fi
\usepackage{bookmark}
\IfFileExists{xurl.sty}{\usepackage{xurl}}{} % add URL line breaks if available
\urlstyle{same}
\hypersetup{
  hidelinks,
  pdfcreator={LaTeX via pandoc}}

\author{}
\date{}

\begin{document}

{
\setcounter{tocdepth}{3}
\tableofcontents
}
\section{Fraktální výpočetní algebra dělitelů nuly --- generativní
operátory, paměť a paralelní
algoritmy}\label{fraktuxe1lnuxed-vuxfdpoux10detnuxed-algebra-dux11blitelux16f-nuly-generativnuxed-operuxe1tory-pamux11bux165-a-paralelnuxed-algoritmy}

\textbf{Samostatný článek z matematiky a informatiky · 25. října 2026}\\
\textbf{Autor:} Vladimír Bíža\\
\textbf{Téma:} Signed algebraické orákulum, implicitní operátory,
úložiště, CPU/GPU a návrh ověřitelných benchmarků.

\subsection{Abstrakt}\label{abstrakt}

Převádíme exaktní informační struktury Paperu A na návrh runtime, v němž
jsou genealogická adresa, incidence, fáze a matematické operátory
rekonstruovatelné bez explicitní materializace celé vyšší algebry.
Výsledkem není univerzální tvrzení, že zero-divisor algebra zrychlí
libovolný program: je to konkrétní systém \textbf{implicitní dědičnost +
explicitní novost + operátorově typovaná sémantika}, který může přímo
sloužit pro adresovatelnou algebraickou paměť, multiresoluční GPU
výpočty, řídké databáze a originální generativní datové indexy.
Přikládáme samostatný Rust oracle, regresní testy a benchmarkový
kontrakt.

\subsection{1. Výpočetní primitivum není uzel stromu, ale strukturální
adresa}\label{vuxfdpoux10detnuxed-primitivum-nenuxed-uzel-stromu-ale-strukturuxe1lnuxed-adresa}

Definujme \(a\in[0,2^k)\), \(\chi=a\oplus b\), kořen \((k_0,\chi_0)\),
ancestry slovo \(u\in\{0,1\}^{k-k_0}\), orientaci a typ operátoru. Čistě
algebraický vztah \(e_a e_b=\varepsilon(a,b)e_{a\oplus b}\) dostaneme
přes rekurzi v čase \(O(k)\), prostoru \(O(k)\) (při iterativním bitovém
provedení \(O(1)\) pomocné paměti). Dvojice assessorů má rozhodnutí
nulové hrany čtyřmi voláními oracle.

\textbf{Kanonický klíč:}
\texttt{CDAddr\ \{\ level,\ a,\ b,\ chi,\ orientation,\ sign,\ ancestry,\ birth\_level,\ relation\_kind\ \}}.
\texttt{chi} nikdy nenahrazuje čtyři adresy; \texttt{relation\_kind}
musí rozlišit ZD od non-ZD, comm, assoc, alt, shared a birth. Tento klíč
je strukturální adresa algebraické identity, nikoli automatické uložení
jejího sémantického obsahu.

\subsection{2. Nulové matice neukládat; aplikovat
je}\label{nulovuxe9-matice-neukluxe1dat-aplikovat-je}

Zděděná incidence \(A_t=A_0\otimes J_m\), \(m=2^t\), umožňuje aplikaci
na rozdělený vektor \(x_{p,u}\) bez jediného alokovaného \(6m\times6m\)
bloku:

\[
s_q=\sum_{v=1}^{m} x_{q,v},\quad r_p=\sum_q (A_0)_{pq}s_q,\quad y_{p,u}=r_p.
\tag{1}
\]

\textbf{Náklad:} obecná dense aplikace \(O((6m)^2)\), generativní
aplikace \(O(6m+6^2)\), paměť \(O(6+m\cdot6)\) pro vstup a výstup,
\(O(6)\) pracovní sumy. V případě nematerializovaného/streamovaného
výstupu lze i výstup konstruovat lazy. Pro \(m=128\) by měla jedna
šestikořenová matice 589824 explicitních prvků; generativní operátor
uchovává pouze 6×6 kořen a velikost rodiny.

\begin{Shaded}
\begin{Highlighting}[]
\KeywordTok{def}\NormalTok{ apply\_inherited(a0, x):}
    \CommentTok{\# x[p][u], p in 0..root{-}1, u in 0..m{-}1}
\NormalTok{    sums }\OperatorTok{=}\NormalTok{ [}\BuiltInTok{sum}\NormalTok{(block) }\ControlFlowTok{for}\NormalTok{ block }\KeywordTok{in}\NormalTok{ x]}
\NormalTok{    out }\OperatorTok{=}\NormalTok{ [}\BuiltInTok{sum}\NormalTok{(a0[p][q] }\OperatorTok{*}\NormalTok{ sums[q] }\ControlFlowTok{for}\NormalTok{ q }\KeywordTok{in} \BuiltInTok{range}\NormalTok{(}\BuiltInTok{len}\NormalTok{(a0)))}
           \ControlFlowTok{for}\NormalTok{ p }\KeywordTok{in} \BuiltInTok{range}\NormalTok{(}\BuiltInTok{len}\NormalTok{(a0))]}
    \ControlFlowTok{return}\NormalTok{ [[out[p]] }\OperatorTok{*} \BuiltInTok{len}\NormalTok{(x[p]) }\ControlFlowTok{for}\NormalTok{ p }\KeywordTok{in} \BuiltInTok{range}\NormalTok{(}\BuiltInTok{len}\NormalTok{(a0))]}
\end{Highlighting}
\end{Shaded}

Pozor: jde o adjacency zděděné komponenty, nikoli o celou signed CD
multiplikaci a nikoli o všechny born komponenty.

\subsection{3. Inherited vs birth: jak reprezentovat skutečnou
novost}\label{inherited-vs-birth-jak-reprezentovat-skuteux10dnou-novost}

Pro každé zdvojení uložíme: (a) kořenový nebo již známý zděděný faktor,
(b) volací pravidlo sign rekurze, (c) \textbf{birth frontier}, který
nese jen nově vzniklou kvalifikovanou vazbu, a (d) volitelné cached
čtení. Kombinatorický budget splňuje \(B_{2N}=E_{2N}-4E_N\). Samotný
počet však nevytváří seznam nových adres: generátor musí nové \(\chi\) a
čtyřadresový oracle skutečně procházet nebo jinak konstruktivně
enumerovat. Výsledný runtime by měl umět vrátit nejen \texttt{is\_zero}
ale i \texttt{why}: provenience kořene, signed predikát, reprezentanta
třídy a úroveň narození.

\subsection{4. Algebraické indexy pro AI a stavové
systémy}\label{algebraickuxe9-indexy-pro-ai-a-stavovuxe9-systuxe9my}

Genealogii používáme jako strukturovaný \textbf{algebraický index},
nikoli jako automaticky naučenou reprezentaci reálných dat. V AI,
distribuovaném zpracování a aplikačních grafech oddělujeme matematický
klíč, reálně zaznamenaný stav a sémantický obsah. Počty a signatury
vztahů se dopočítávají algoritmem, skutečné uživatelské události a
parametry modelu se ukládají.

Výpočetní pipeline uchovává \texttt{relation\_kind} pro \texttt{ZD},
\texttt{non-ZD}, \texttt{comm}, \texttt{assoc}, \texttt{alt} a
\texttt{birth}. Nový vstup může upravit existující větev nebo přidat
samostatnou \texttt{birth} větev; matematické zdvojení samo o sobě
nevytváří zkušenost, fakta ani naučené váhy. Operátory trojných součinů
musí zachovat závorkování. Pro pozorovatelnost používáme hodnost
Krylovovy matice a požadujeme další sondu právě tehdy, když chybí
rozměry readoutu.

\subsection{5. Paměťová pyramida: VRAM / RAM /
NVMe}\label{pamux11bux165ovuxe1-pyramida-vram-ram-nvme}

\begin{itemize}
\tightlist
\item
  \textbf{VRAM:} aktivní amplitudy, hot blocky, GPU batch pro
  \(\varepsilon\), Krylovovy báze, lokální operátory a jejich dočasné
  projekce.
\item
  \textbf{RAM:} ancestry index, hot/warm factor cache, recent birth
  frontier a prefetch plány.
\item
  \textbf{NVMe:} committed stav, provenience, immutable kořeny,
  transakční log, nové sémantické objekty a kvalifikované ZD frontier.
\end{itemize}

Vrstvy přesouvají \textbf{data a aktivní support}, nikoli matematickou
definici pléna. Čisté \texttt{K\_\{2048\}} neimplikuje alokaci
\(2048^2\) tabulky. Strukturované operátory mohou být implicitní;
samotný živý stav však pro hustých 2048 amplitud stále vyžaduje
příslušné amplitudy.

\subsection{6. GPU paralelismus --- rozdělovat podle algebraické
práce}\label{gpu-paralelismus-rozdux11blovat-podle-algebraickuxe9-pruxe1ce}

Efektivní kernelová hranice je typicky
\texttt{chi\ ×\ ancestry-prefix\ ×\ relation\_kind\ ×\ active-support}.
Pro zděděné segmenty použijeme redukce bloků a broadcast podle (1). Pro
born segmenty frontier spustí řídký generátor nebo sparse gather. Pro
výpočet signed násobení je přirozené seskupení XOR a bitové rekurze;
paměťový layout je \texttt{struct-of-arrays} pro koalescentní GPU fetch.
Neslučujeme různé vztahové kanály do jediného booleovského grafu.

\begin{Shaded}
\begin{Highlighting}[]
\NormalTok{stage A: select active (root, chi, ancestry, relation{-}kind)}
\NormalTok{stage B: classify inherited factor / born frontier}
\NormalTok{stage C: batch sign oracle + sparse algebraic operators}
\NormalTok{stage D: reduce partial responses deterministically}
\NormalTok{stage E: collect closure residual + persist qualified delta}
\end{Highlighting}
\end{Shaded}

Synchronizace je nutná tam, kde \texttt{assoc} nebo \texttt{alt} propojí
odlišné \(\chi\)-bloky. Deterministické sloučení musí uchovat orientaci,
sign a závorkování.

\subsection{7. Databáze a obecné
stromy}\label{databuxe1ze-a-obecnuxe9-stromy}

Jde o dvě různé úrovně implementace: \texttt{source\ algebra\ index} pro
matematickou sémantiku a \texttt{inspired\ hierarchical\ index} pro
běžná data. U obecného dokumentového či aplikačního stromu
nepředstíráme, že jeho podniková logika vznikla z CD násobení; přebíráme
mechanismy strukturálního sdílení a delta reprezentace.

\textbf{Návrh schématu PostgreSQL:}

\begin{Shaded}
\begin{Highlighting}[]
\KeywordTok{CREATE} \KeywordTok{TABLE}\NormalTok{ zd\_root (}
\NormalTok{  root\_id BIGINT }\KeywordTok{PRIMARY} \KeywordTok{KEY}\NormalTok{,}
  \KeywordTok{level} \DataTypeTok{INT} \KeywordTok{NOT} \KeywordTok{NULL}\NormalTok{,}
\NormalTok{  chi BIGINT }\KeywordTok{NOT} \KeywordTok{NULL}\NormalTok{,}
\NormalTok{  scheme\_version }\DataTypeTok{INT} \KeywordTok{NOT} \KeywordTok{NULL}\NormalTok{,}
\NormalTok{  immutable\_hash BYTEA }\KeywordTok{NOT} \KeywordTok{NULL}
\NormalTok{);}
\KeywordTok{CREATE} \KeywordTok{TABLE}\NormalTok{ zd\_delta (}
\NormalTok{  root\_id BIGINT }\KeywordTok{REFERENCES}\NormalTok{ zd\_root(root\_id),}
\NormalTok{  ancestry BYTEA }\KeywordTok{NOT} \KeywordTok{NULL}\NormalTok{,}
\NormalTok{  birth\_level }\DataTypeTok{INT} \KeywordTok{NOT} \KeywordTok{NULL}\NormalTok{,}
\NormalTok{  a BIGINT }\KeywordTok{NOT} \KeywordTok{NULL}\NormalTok{, b BIGINT }\KeywordTok{NOT} \KeywordTok{NULL}\NormalTok{,}
\NormalTok{  c BIGINT }\KeywordTok{NOT} \KeywordTok{NULL}\NormalTok{, d BIGINT }\KeywordTok{NOT} \KeywordTok{NULL}\NormalTok{,}
\NormalTok{  relation\_kind }\DataTypeTok{SMALLINT} \KeywordTok{NOT} \KeywordTok{NULL}\NormalTok{,}
\NormalTok{  orientation }\DataTypeTok{SMALLINT} \KeywordTok{NOT} \KeywordTok{NULL}\NormalTok{,}
\NormalTok{  payload\_ref BYTEA,}
  \KeywordTok{PRIMARY} \KeywordTok{KEY}\NormalTok{(root\_id, ancestry, a,b,c,d,relation\_kind)}
\NormalTok{);}
\KeywordTok{CREATE} \KeywordTok{INDEX}\NormalTok{ zd\_delta\_birth\_idx }\KeywordTok{ON}\NormalTok{ zd\_delta(birth\_level, relation\_kind);}
\end{Highlighting}
\end{Shaded}

Store only qualified novelty and actual semantic payload; inherited
mathematically determined structure is calculated on reads. Transakční
log zaručuje committed stav. Celočíselné adresy nad \texttt{BIGINT}
vyžadují širší bitset/byte encoding.

\subsection{8. Higher-π jako výpočetní orákulum a optimalizační
metrika}\label{higher-ux3c0-jako-vuxfdpoux10detnuxed-oruxe1kulum-a-optimalizaux10dnuxed-metrika}

Algebraické invarianty \texttt{V\_N}, \texttt{Phi\_N}, \texttt{Q\_N},
\texttt{Q\_phase}, \texttt{C\_N} poskytují přesné jednotkové testy bez
Monte Carla. \texttt{Q\_N} kvantifikuje průměrný čtvercový defekt
násobení normovaných zdrojových prvků; je to analytická charakteristika
jednotlivé CD vrstvy, kterou lze použít pro numerické testování signed
kernelů, kontrolu implementace a výběr reprezentace podle požadované
přesnosti.

\begin{Shaded}
\begin{Highlighting}[]
\ImportTok{from}\NormalTok{ fractions }\ImportTok{import}\NormalTok{ Fraction}
\ImportTok{from}\NormalTok{ math }\ImportTok{import}\NormalTok{ comb}

\KeywordTok{def}\NormalTok{ edges(n):}
    \ControlFlowTok{return} \DecValTok{0} \ControlFlowTok{if}\NormalTok{ n }\OperatorTok{\textless{}} \DecValTok{16} \ControlFlowTok{else}\NormalTok{ (n}\OperatorTok{{-}}\DecValTok{2}\NormalTok{)}\OperatorTok{*}\NormalTok{(n}\OperatorTok{{-}}\DecValTok{4}\NormalTok{)}\OperatorTok{*}\NormalTok{(n}\OperatorTok{{-}}\DecValTok{8}\NormalTok{)}\OperatorTok{//}\DecValTok{16}

\KeywordTok{def}\NormalTok{ q(n):}
    \ControlFlowTok{return}\NormalTok{ Fraction(}\DecValTok{32}\OperatorTok{*}\NormalTok{edges(n), n}\OperatorTok{*}\NormalTok{n}\OperatorTok{*}\NormalTok{(n}\OperatorTok{+}\DecValTok{2}\NormalTok{)}\OperatorTok{**}\DecValTok{2}\NormalTok{)}

\KeywordTok{def}\NormalTok{ q\_phase(n):}
    \ControlFlowTok{return}\NormalTok{ Fraction(}\DecValTok{32}\OperatorTok{*}\NormalTok{edges(n), (n}\OperatorTok{{-}}\DecValTok{1}\NormalTok{)}\OperatorTok{**}\DecValTok{2}\OperatorTok{*}\NormalTok{(n}\OperatorTok{+}\DecValTok{1}\NormalTok{)}\OperatorTok{**}\DecValTok{2}\NormalTok{)}

\KeywordTok{def}\NormalTok{ compound(n):}
    \ControlFlowTok{return}\NormalTok{ Fraction(comb(n,n}\OperatorTok{//}\DecValTok{2}\NormalTok{), }\DecValTok{1}\NormalTok{) }\OperatorTok{*}\NormalTok{ (}\DecValTok{1}\OperatorTok{+}\NormalTok{q(n))}\OperatorTok{**}\DecValTok{2} \OperatorTok{/}\NormalTok{ (}\DecValTok{1}\OperatorTok{+}\NormalTok{q(}\DecValTok{2}\OperatorTok{*}\NormalTok{n))}

\ControlFlowTok{assert}\NormalTok{ q(}\DecValTok{16}\NormalTok{) }\OperatorTok{==}\NormalTok{ Fraction(}\DecValTok{7}\NormalTok{,}\DecValTok{216}\NormalTok{)}
\ControlFlowTok{assert}\NormalTok{ q(}\DecValTok{32}\NormalTok{) }\OperatorTok{==}\NormalTok{ Fraction(}\DecValTok{315}\NormalTok{,}\DecValTok{9248}\NormalTok{)}
\ControlFlowTok{assert}\NormalTok{ compound(}\DecValTok{8}\NormalTok{) }\OperatorTok{==}\NormalTok{ Fraction(}\DecValTok{15120}\NormalTok{,}\DecValTok{223}\NormalTok{)}
\end{Highlighting}
\end{Shaded}

\subsection{9. Reprodukovatelnost a výkonnostní
kontrakt}\label{reprodukovatelnost-a-vuxfdkonnostnuxed-kontrakt}

\textbf{Dodané:} \texttt{code/src/main.rs}: rekurzivní signed-CD oracle,
signed nulový predikát, explicitní svědek a Rust unit testy pro
\texttt{K16}, \texttt{K32}, four-lift. Výpočty K16/K32/K64 a svědka byly
nezávisle numericky ověřeny Python ekvivalentem. Rust kompilátor nebyl v
tvorbě balíku dostupný; zdroj je připraven ke spuštění příkazy
\texttt{cargo\ test} a \texttt{cargo\ run\ -\/-release\ -\/-\ 32}.

\textbf{Benchmark vůči CSR:} pro stejné úlohy změřit (i) generování
\(E_N\), (ii) \texttt{apply(A,x)}, (iii) random ancestor lookup, (iv)
born-frontier inkrement, (v) persist+reload, (vi) Krylov probes, (vii)
celé end-to-end task. Uvádět peak RSS/VRAM, propustnost, p50/p99,
přesnost a bajty na kvalifikovanou vazbu. Faktorované zděděné bloky mají
algoritmickou výhodu; celkové produkční zrychlení se musí spočítat na
konkrétní pracovní zátěži.

\subsection{10. První výpočetní
syntéza}\label{prvnuxed-vuxfdpoux10detnuxed-syntuxe9za}

ZD genealogie přináší výpočetní princip, ve kterém je
\textbf{matematicky vynucený vztah algoritmem, nikoli položkou
databáze}. Higher-π poskytuje přesné kontrolní invarianty téhož
algebraického zdroje; born-frontier zachovává skutečnou novost a
multiplex udržuje nonalternativní informaci. Výsledkem je uskutečnitelný
základ pro rozlehlé implicitní paměťové a výpočetní světy, včetně
stavové a AI systémy --- ne přilepený ZD graf, ale řízená generativní
operátorová topologie.

\subsection{11. Úplný typovaný kontrakt a algebraické
invarianty}\label{uxfaplnuxfd-typovanuxfd-kontrakt-a-algebraickuxe9-invarianty}

Dělíme identitu algebraického objektu, jeho sémantický stav a fyzickou
adresu uložení. \texttt{AlgebraKey} obsahuje \texttt{tower\_version},
\texttt{level}, XOR barvu, dvě kanonizované dvojice bázových adres,
relativní sign a \texttt{relation\_kind}; \texttt{StorageKey} navíc
kóduje \texttt{root\_id}, ancestry, epochu a verzi schématu.
\texttt{SemanticPayload} je nezávislá content-addressed data struktura.
Signed oracle nikdy nevytváří historickou zkušenost uživatele, pouze
ověřuje algebraickou povolenost vazby.

Závorkování není metadata na okraji: strom výrazu
\texttt{Mul(Mul(x,y),z)} se musí rozlišovat od \texttt{Mul(x,Mul(y,z))}.
Při přesunu mezi CUDA a CPU musí referenční exact oracle ověřit znak,
podporu výsledku i asociátor/alternátor. Pro floating-point amplitudy
testujeme zvlášť numerickou odchylku; algebraický index sám čísla
nezaokrouhluje.

\subsection{12. Výpočetní jádra a nákladový
model}\label{vuxfdpoux10detnuxed-juxe1dra-a-nuxe1kladovuxfd-model}

Nechť root incidence má \(r\) uzlů, \(e\) nenulových hran a rodina má
\(m=2^t\) následníků na root. Plný explicitní blok obsahuje \(r^2m^2\)
prvků. Generativní \texttt{apply} má cenu \(\Theta(rm+e)\) pro redukci
vstupu, sparse root aplikaci a materializaci výstupu; pracovní paměť
kromě vstupu/výstupu \(O(r)\). Pro náhodné čtení jednotlivé hrany stačí
klasifikace root dvojice a ancestry, bez alokace celé matice.

Algoritmus pro řídký birthday frontier odděluje tři kroky:
deterministickou enumeraci kandidátních assessorů, rozhodnutí \(W=+1\) a
deduplikaci výsledných klíčů. \texttt{birth\_count} samotný nenahrazuje
konstruktivní enumerátor. Při paměťové kompresi
\(\Theta(m^2)\to\Theta(1)\) jde o \textbf{reprezentaci jedné pravidelné
komponenty}, nikoli o univerzální kompresní poměr všech dat či
libovolného grafu.

\subsection{13. Paralelní determinismus a
persistence}\label{paralelnuxed-determinismus-a-persistence}

Work units seskupujeme podle
\texttt{(level,\ chi,\ root,\ ancestry\_prefix,\ channel)} a přidělujeme
jim monotonický \texttt{epoch}. Každá vrstva produkuje immutable
\texttt{Delta\{epoch,\ source\_hash,\ op,\ payload\_hash\}}. Redukce
musí mít jednoznačné pořadí při floating-point aritmetice;
integer/symbolic výsledky mohou používat komutativní agregaci jen tam,
kde to dovoluje algebraická operace. \texttt{commit} probíhá až po
validaci všech rodičovských adres, externí sémantické payloady zůstávají
obsahově adresované.

PostgreSQL transakčně ukládá \emph{zápisy}, nikoli generativní pravidla
přepočítávaná jako data. NVMe snapshot obsahuje epochu, verzi sign
rekurze a root hash; cold reload provede recalculation oracle a ověří
kontrolní součty. Multiplexní dotaz vrací \texttt{ZD}, \texttt{non-ZD},
\texttt{comm}, \texttt{assoc}, \texttt{alt}, \texttt{birth} odděleně, a
pokud sondy neobsáhnou celý Krylovův prostor, uvádí přesný rozměr
nepozorovatelného komplementu.

\subsection{14. Výpočetní experiment a regresní
protokol}\label{vuxfdpoux10detnuxed-experiment-a-regresnuxed-protokol}

Dodané \texttt{code/benchmark.py} provádí přesně stejný \texttt{apply} s
explicitní referenční maticí a faktorizovaným blokem. Výstupy se
kontrolují na rovnost a ukládají se jako CSV s typem experimentu,
rozměrem, dobou a odhadovaným objemem materializovaných dat. Nepoužíváme
naměřené doby jedné konkrétní platformy jako univerzální hardwarový
závěr. \texttt{code/verify.py} certifikuje základní číselné identity;
\texttt{code/src/main.rs} reprodukuje signed oracle nezávislou jazykovou
implementací.

Učení a změnu kompetencí systému vyhodnocujeme třemi nezávislými testy:
(a) zachování exaktní algebraické identity po refinementu, (b) správnost
persistovaných pracovních a trvalých stavů, (c) změna skutečných
kognitivních kompetencí v teacher-removed testu. Pouhé snížení VRAM není
doklad učení a úspěšné algebraické testy nejsou dokladem obecné
inteligence.

\subsection{15. Závěr}\label{zuxe1vux11br}

Jednotný přínos je konkrétní: signed CD definice vytváří adresovatelné
vztahy, genealogie umožňuje faktorizované implicitní výpočty, birth
frontier zachycuje nové členy a normově-momentové invarianty jsou
přesnými regresními kontrolami. Tento návrh je realizovatelný bez
materializace celé CD tabulky; strukturální úspory jsou vyjádřeny
algoritmickými náklady a doprovodným benchmarkem.

\subsection{Literatura}\label{literatura}

{[}1{]} V. Bíža, \emph{Fraktální výpočetní algebra dělitelů nuly ---
signed Cayley--Dicksonův informační les a vyšší π}, doprovodný Paper
A.\\
{[}2{]} G. Moreno, \emph{Constructing zero divisors in the higher
dimensional Cayley--Dickson algebras}, arXiv:math/0512517.\\
{[}3{]} T. A. Davis, \emph{Algorithm 1000: SuiteSparse:GraphBLAS}, ACM
Transactions on Mathematical Software.

\end{document}
